\documentclass[9pt,letterpaper]{extarticle}
\usepackage[margin=0.58in]{geometry}
\usepackage[T1]{fontenc}
\usepackage{lmodern}
\usepackage{amsmath,amssymb,mathtools,booktabs,array,enumitem,microtype}
\usepackage[colorlinks=true,linkcolor=blue,urlcolor=blue]{hyperref}
\setlength{\parindent}{0pt}
\setlength{\parskip}{3.5pt}
\setlist[itemize]{leftmargin=1.35em,itemsep=1.5pt,topsep=2pt,parsep=0pt}
\newcommand{\Z}{\mathbb Z}
\newcommand{\C}{\mathbb C}
\newcommand{\Pone}{\mathbb P^1}
\newcommand{\M}{\mathcal M}
\pagestyle{plain}
\begin{document}
\begin{center}
{\Large\bfseries The $S^6$ Complex-Threefold Construction}\\[-1pt]
{\large A two-page mathematical status after proof reconstruction}\\[2pt]
{\small Dated corrective audit --- 25 August 2026}
\end{center}

\textbf{Current disposition.}
The earlier audit's topology rejection is withdrawn as stated.  The manuscript contains two mathematically related specialization routes, and the first audit treated a bridge needed by the explicit collapse route as though it were needed by both.  The non-collapse route through intrinsic nearby cycles, Leray transgression, and based van Kampen survives the present line-by-line reconstruction.  The resulting review is provisionally favorable, with focused checks still requested for the global analytic construction and the external theorems invoked in nearby-cycle, Gysin, duality, and base-change arguments.

\textbf{Construction and arithmetic datum.}
The manuscript constructs a compact complex threefold
\[
 f:X\longrightarrow \Pone
\]
with generic fibre a two-dimensional complex torus, one reduced non-normal normal-crossings fibre $W$, and two multiple fibres of multiplicities $3$ and $4$.  On the rank-four lattice, the exact monodromy calculation gives
\[
T_1^3=T_2^4=I,\qquad T_0=(T_1T_2)^{-1}=I+N,\qquad N^2=0,\quad \operatorname{rk}N=2.
\]
The logarithmic twists satisfy
\[
(\ell_0,\ell_1,\ell_2)=(0,1,-1),\qquad
p=12\ell_0-4\ell_1-3\ell_2=-1.
\]
That same integer controls both the fundamental-group presentation and the three Leray transgressions.

\textbf{The morphism-level correction.}
The explicit cusp-collapse map $c:F\to W$ supplies
\[
c^*:H^q(W;\Z)\longrightarrow H^q(F;\Z),
\]
whereas the nearby-cycle construction supplies the intrinsic stalk map
\[
\operatorname{sp}^q:H^q(W;\Z)\longrightarrow H^q(F;\Z)^{T_0}.
\]
The manuscript identifies them when the collapse construction is used, but Section~7.7 requires only $\operatorname{sp}^q$.  Appendix~B proves this map independently of $c$ by the local models $t=z_0$, $t=z_0z_1$, $t=z_0z_1z_2$, the associated nearby-cycle sheaves, the conductor/normalization calculation on $W$, hypercohomology rank bookkeeping, a degree-three support/Gysin argument, and integral saturation.  The resulting maps are
\[
\operatorname{sp}^q:H^q(W;\Z)\xrightarrow{\sim}H^q(F;\Z)^{T_0},\qquad 0\le q\le4,
\]
with invariant ranks $(1,2,4,2,1)$, independently reproduced by exact exterior-power and Smith-normal-form calculations.

\textbf{Leray and fundamental group.}
The central stalk above and the finite-fibre indices give
\[
H^0(B,R^1f_*\Z)=12\Z\gamma,\quad
H^0(B,R^2f_*\Z)=2\Z q,\quad
H^0(B,R^3f_*\Z)=2\Z\gamma uw.
\]
Local lift discrepancies around the three based meridians produce
\[
d_2^{0,1}(12\gamma)=\pm p\omega,
\]
and multiplicativity gives, with the same sign,
\[
d_2^{0,2}(2q)=\pm p\,\omega\bar\delta,
\qquad
d_2^{0,3}(2\gamma uw)=\pm p\,\omega[\gamma\delta].
\]
Since $p=-1$, all three maps are units, so $H^1(X;\Z)=H^2(X;\Z)=H^3(X;\Z)=0$; duality and universal coefficients give the integral homology of $S^6$.

The based overlap calculation uses explicit deck transformations and gives
\[
\pi_1(X)=\langle c,x,y\mid c\text{ central},\;x^3=c^{\ell_1},\;y^4=c^{\ell_2},\;xy=c^{\ell_0}\rangle.
\]
Its relation determinant is $-p$, its actual Smith form is $\operatorname{diag}(1,1,1)$, and the nonabelian elimination yields a cyclic group of order $|p|$; hence $\pi_1(X)=1$.  Hurewicz, Whitehead, and $\Theta_6=0$ then identify the underlying smooth manifold with $S^6$.

\newpage
\textbf{Independent algebraic-dimension result.}
For an integral bivector $E$, the fibrewise $(1,1)$ condition is the vanishing of
\[
P_E=E_{\gamma u}-E_{\gamma w}\tau-E_{\gamma\delta}\mu
 +6E_{uw}\mu+E_{u\delta}\beta+E_{w\delta}(6\mu^2-\tau\beta).
\]
The transformation laws under the two generators prove linear independence of
\[
1,\ \tau,\ \mu,\ \beta,\ 6\mu^2-\tau\beta.
\]
Thus the generic integral kernel is generated by
\[
\eta=u\wedge w+6\gamma\wedge\delta,\qquad \eta^2=12.
\]
The associated Hermitian form has determinant $\frac23\operatorname{Im}\tau\,D<0$ and signature $(1,1)$.  The generic fibre therefore has N\'eron--Severi group $\Z\eta$, while the torus algebraic-reduction alternatives of dimensions $1$ and $2$ are excluded.  A meromorphic function on $X$ is constant along a very general fibre because the analytic locus where its vertical differential vanishes contains those fibres and hence all of $X$.  It descends to the base, giving
\[
\M(X)=f^*\C(t),\qquad a(X)=1,
\]
and $f$ is the algebraic reduction.  This proof does not use Section~7 or the CDP comparison.

\textbf{Section 10 mechanism.}
For a reduced normal-crossings non-normal fibre, the normalization/conductor sequence is exact in the local models $xy=0$ and $xyz=0$.  The ambient differential of the defining equation has coefficients
\[
(y,x,0)\quad\text{or}\quad(yz,xz,xy),
\]
so it vanishes on the conductor and glues to a nonzero ambient one-form whose intrinsic image is conductor-supported torsion.  This verifies the local mechanism by which the ambient differential can evade a conormal reduction.  Its global deployment uses the previously established topology and Picard information together with analytic Serre--Grothendieck duality and top-degree base change; those external inputs remain explicit review targets.

\textbf{What remains to be checked.}
\begin{itemize}
  \item The explicit cusp-collapse map and the first Mayer--Vietoris route remain open as separate claims: continuity across every toric stratum, quotient descent, and the resulting integral collapse matrices have not been fully rederived at chain level.  They are redundant for the non-collapse recognition route but still matter for the manuscript's claimed first topology proof and geometric description of $W$.
  \item Sections~3--6 have been read for coherence, and their finite transformation/freeness subclaims pass exact checks.  A publication-level audit should still verify every proper-discontinuity, Hausdorffness, compactness, and holomorphic-gluing hypothesis lemma by lemma.
  \item Appendix~B and Section~10 use standard external formalism.  The next high-value task is a theorem-by-theorem source audit recording the exact hypotheses, naturality statements, signs, and integral saturation conclusions used here.
\end{itemize}

\textbf{Provenance.}
The candidate manuscript supplies the original construction.  The Product/Codex workbench identified the independent Leray and direct meromorphic-descent routes and repaired two transcription-level mistakes in its own reconstruction.  The present integration treated that progress record as a roadmap, returned to the canonical 108-page source, reproduced the finite algebra, and reclassified each route separately.  No claim here relies on inaccessible live workbench files.

\vfill
{\footnotesize Evidence pointers: revised proof audit; revised theorem and risk ledgers; 17-check integration harness; prior 38-check core harness; revised 108-page annotation overlay.}
\end{document}
